%!TEX root = ../example.tex
\chapter{Introducción}
\label{ch:introduccion}

    El presente documento contiene la especificación del entregable \varCveEntregable:\varEntregable\ del Proyecto \varProyecto, correspondiente a la plataforma web funcional y su documentación técnica respectiva. En este reporte se detallan los avances en el análisis y diseño del sistema orientado a la seguridad digital institucional.

    \section{Intención del documento}

        En este documento refleja el resultado de los avances realizados en el sistema \varCveSistema-\varSistema\ y tiene como objetivo fundamental establecer la arquitectura, los módulos de interacción y las reglas de negocio necesarias para el desarrollo de un canal digital seguro de denuncia y prevención.

        Va dirigido a los directores del Trabajo Terminal, el M. en C. Ulises Vélez Saldaña y la Dra. Jessie Paulina Guzmán Flores, así como a los sínodos encargados de la evaluación técnica del proyecto y a los desarrolladores involucrados en la implementación de los módulos de Backend y Frontend.

    \section{Estructura del documento}

        El presente documento se encuentra organizado en tres partes principales, cada una dividida a su vez en capítulos, como se muestra a continuación:

        \begin{itemize}
            \item El capítulo \ref{ch:nomenclatura} presenta la terminología y modelado base utilizado para la estructura de la información.
            \item Parte 1: Modelo de Negocio
                \begin{itemize}
                    \item El capítulo \ref{ch:glosario} presenta los términos técnicos y del negocio necesarios para comprender el dominio de la denuncia segura y riesgos digitales.
                    \item El capítulo \ref{ch:modInfo} detalla las entidades del sistema como Alumno, Domicilio y Proceso de Admisión.
                    \item El capítulo \ref{ch:modReglas} define las reglas de negocio que rigen el comportamiento del sistema.
                \end{itemize}
            \item Parte 2: Modelo Dinámico e Interacción
                \begin{itemize}
                    \item El capítulo \ref{ch:modDinamico} describe los casos de uso, como el llenado de solicitudes y gestión de información.
                    \item El capítulo \ref{ch:modAlcance} define los actores del sistema, incluyendo al Gerente de Ventas y otros roles administrativos.
                    \item El capítulo de interfaces presenta el diseño visual y flujo de interacción del usuario.
                \end{itemize}
        \end{itemize}
    
\chapter{Nomenclatura}
\label{ch:nomenclatura}
    
    La información del presente documento se encuentra estructurada mediante estándares de ingeniería de software para la documentación de requerimientos, utilizando diagramas UML y tablas descriptivas para garantizar la claridad en la especificación de componentes.

    \section{Modelado del negocio}
    
        Para modelar el negocio, se presenta el glosario de términos, el modelado de la estructura de la información que tendrá el sistema, las reglas de negocio que rigen las operaciones y máquinas de estados para modelar los ciclos de gestión en el proyecto \varProyecto.

        \subsection{Modelado de estructura de datos}
        \label{sec:ModEstructuraDatos}
        Se utiliza el modelado de entidades para definir los atributos y relaciones de los datos persistentes en el sistema, asegurando la integridad de la información del usuario y los reportes.
        
        \subsection{Reglas de Negocio}
        \label{sec:ModReglasNegocio}
        Se especifican las restricciones y políticas operativas, tales como la obligatoriedad de campos y el formato de los datos ingresados, para mantener la coherencia del sistema.
        
        \subsection{Maquinas de Estados}
        \label{sec:ModMaquinas}
        Se emplean para describir el ciclo de vida de los procesos principales, como el cambio de estatus de una denuncia desde su registro hasta su resolución final.
    
    \clearpage
    \section{Modelado de casos de uso}
    Se describen las funciones del sistema desde la perspectiva del usuario, detallando trayectorias principales y alternativas para acciones como el llenado de solicitudes.

    \section{Modelado de interfaces}
    Presenta el diseño de las pantallas del sistema, facilitando la comprensión del flujo de navegación y la interacción del actor con los módulos de inscripción y consulta.

    \section{Modelado de mensajes}
    Detalla los avisos de operación exitosa, errores y confirmaciones que el sistema emite para guiar al usuario durante su navegación.